\section{The model and its transition density}
\label{sec:model}

\subsection{The stochastic differential equation}

The library implements the arithmetic jump-diffusion
\begin{equation}
  \dd X_t \;=\; \mu \dd t \;+\; \sigma \dd W_t \;+\; J_t \dd N_t,
  \label{eq:sde}
\end{equation}
where $N_t$ is a Poisson process with intensity $\lambda$, independent of
$W_t$, and the jump sizes $J_t$ are i.i.d.\ draws from
$f_J(\cdot\,;\theta_J)$, independent of both. Because the coefficients are
constant, \eqref{eq:sde} integrates in closed form:
\begin{equation}
  X_t \;=\; X_0 \;+\; \mu t \;+\; \sigma W_t \;+\; \sum_{k=1}^{N_t} J_k .
  \label{eq:solution}
\end{equation}
The arithmetic specification is intended for \emph{log-prices}: if $S_t$
is the price and $X_t = \log S_t$, then \eqref{eq:sde} is the standard
Merton model written on returns, and the increments
$\Delta x_i$ are the log-returns the user hands to the estimator.

\begin{remark}
Equation \eqref{eq:solution} is exact, not an Euler approximation ---
constant coefficients make the integral trivial. Every approximation in
the library enters later, through the treatment of the jump count per
observation interval, not through time discretization of the diffusion.
\end{remark}

\subsection{Exact transition density: the Poisson mixture}

Over an interval of length $\Delta t$, the increment implied by
\eqref{eq:solution} is
\begin{equation}
  \Delta X \;=\; \mu\,\Delta t \;+\; \sigma \sqrt{\Delta t}\, Z
  \;+\; \sum_{k=1}^{K} J_k,
  \qquad Z \sim \mathcal N(0,1),\quad
  K \sim \mathrm{Poisson}(\lambda \Delta t),
  \label{eq:increment}
\end{equation}
with $Z$, $K$ and the $J_k$ mutually independent. Conditioning on
$K = k$ and convolving gives the exact transition density as a Poisson
mixture:
\begin{equation}
  g(\Delta x)
  \;=\;
  \sum_{k=0}^{\infty}
  e^{-\lambda \Delta t}\,\frac{(\lambda \Delta t)^k}{k!}
  \left( \phi_{\mu \Delta t,\, \sigma^2 \Delta t} * f_J^{*k} \right)
  (\Delta x),
  \label{eq:poisson-mixture}
\end{equation}
where $f_J^{*k}$ is the $k$-fold convolution of the jump density
($f_J^{*0} = \delta_0$). Each additional term requires one more
convolution, and for most jump laws none of them is available in closed
form.

\subsection{The Bernoulli approximation}
\label{sec:bernoulli}

\begin{decision}[At most one jump per interval]
\label{dec:bernoulli}
For daily data, $\lambda \Delta t$ is small: at $\lambda \Delta t = p$,
the probability of two or more jumps in one interval is
\[
  \Prob(K \ge 2) \;=\; 1 - e^{-p}(1 + p) \;=\; \tfrac{p^2}{2} + O(p^3),
\]
i.e.\ about $0.005$ even at the fairly extreme $p = 0.1$. The library
therefore truncates \eqref{eq:poisson-mixture} at $k \le 1$ and replaces
the Poisson count by a Bernoulli indicator $B \sim \mathrm{Bernoulli}(p)$
with $p = \lambda \Delta t$:
\begin{equation}
  \Delta X \;\approx\; \mu\,\Delta t + \sigma \sqrt{\Delta t}\, Z + B\,J .
  \label{eq:bernoulli-increment}
\end{equation}
This is the same approximation used in \citet{Ospina2009} and
\citet{Ardia2011}. It buys a two-component mixture density (one
convolution instead of infinitely many) at a bias cost of order
$(\lambda\Delta t)^2$; quantifying that bias precisely across the
$\lambda \Delta t$ range is part of the research agenda of
Section~\ref{sec:closing}.
\end{decision}

Under \eqref{eq:bernoulli-increment}, the transition density becomes the
two-component mixture at the heart of the library,
\begin{equation}
  \boxed{\;
  g_{\thetavec}(\Delta x)
  \;=\;
  (1 - p)\,\phi\!\left(\Delta x;\ \mu\Delta t,\ \sigma^2 \Delta t\right)
  \;+\;
  p \left( \phi_{\mu\Delta t,\,\sigma^2\Delta t} * f_J \right)(\Delta x).
  \;}
  \label{eq:mixture}
\end{equation}
Interpreting \eqref{eq:mixture} reads exactly like the data-generating
story: with probability $1-p$ the day is ``ordinary'' and the increment is
Gaussian; with probability $p$ a jump occurred and the increment is the sum
of the Gaussian part and one draw from the jump law --- hence the
convolution.

\subsection{The code: simulation}

\code{JumpDiffusionModel.simulate} generates paths directly from
\eqref{eq:bernoulli-increment}: a vector of standard normal innovations, a
vector of Bernoulli jump indicators, and a vector of jump sizes drawn from
the jump distribution, combined step by step.

\begin{lstlisting}[caption={Path simulation in \code{models/jump\_diffusion.py} (excerpt).}]
dt = T / n_steps
brownian_innovations = np.random.normal(0, 1, n_steps)
jump_indicators = np.random.binomial(1, self.parameters["jump_prob"], n_steps)
jump_sizes = self.jump_distribution.rvs(self._jump_params(), size=n_steps)

for i in range(n_steps):
    drift = self.parameters["mu"] * dt
    sigma_term = self.parameters["sigma"] * np.sqrt(dt)
    diffusion = sigma_term * brownian_innovations[i]
    jump = jump_indicators[i] * jump_sizes[i]
    path[i + 1] = path[i] + drift + diffusion + jump
\end{lstlisting}

Note that a jump size is drawn for every step but only \emph{applied} when
the indicator fires; this keeps the random-number stream aligned across
parameter values, which is convenient for variance-reduction and
reproducibility in Monte Carlo work.

\subsection{The code: the mixture density}

The model's \code{\_mixture\_density} method evaluates \eqref{eq:mixture}.
The only branch is \emph{how} the convolution term is computed: the jump
distribution's \code{diffusion\_convolved\_pdf} is first asked for a
closed form (Section~\ref{sec:distributions}); if it returns \code{None},
the generic FFT fallback of Section~\ref{sec:fft} takes over.

\begin{lstlisting}[caption={The mixture density \eqref{eq:mixture} in \code{models/jump\_diffusion.py}.}]
def _mixture_density(self, x: np.ndarray, dt: float) -> np.ndarray:
    mu = self.parameters["mu"]
    sigma = self.parameters["sigma"]
    p = self.parameters["jump_prob"]

    diffusion_mean = mu * dt
    diffusion_std = sigma * np.sqrt(dt)
    diffusion_density = norm.pdf(x, loc=diffusion_mean, scale=diffusion_std)

    # Jump + diffusion component: use a closed form if the jump
    # distribution has one, otherwise fall back to FFT convolution.
    jump_params = self._jump_params()
    jump_diffusion_density = self.jump_distribution.diffusion_convolved_pdf(
        x, jump_params, diffusion_mean, diffusion_std
    )
    if jump_diffusion_density is None:
        jump_diffusion_density = self.jump_distribution.fft_convolved_pdf(
            x, jump_params, diffusion_mean, diffusion_std
        )

    return (1 - p) * diffusion_density + p * jump_diffusion_density
\end{lstlisting}

The log-likelihood of the observed increments is then simply
\begin{equation}
  \loglik(\thetavec)
  \;=\;
  \sum_{i=1}^{n} \log g_{\thetavec}(\Delta x_i),
  \label{eq:loglik}
\end{equation}
computed in the model's \code{log\_likelihood} method with one numerical
guard: densities are floored at $10^{-300}$ before taking logs, so a
single increment falling in a region where the numerical density
underflows to zero degrades the likelihood instead of destroying it with
$\log 0 = -\infty$.

\begin{lstlisting}[caption={Log-likelihood \eqref{eq:loglik} with underflow guard.}]
def log_likelihood(self, data: np.ndarray, dt: float) -> float:
    increments = data if data.ndim == 1 else np.diff(data)
    densities = self._mixture_density(increments, dt)
    densities = np.maximum(densities, 1e-300)  # Numerical stability
    return np.sum(np.log(densities))
\end{lstlisting}
